\documentclass[11pt,a4paper]{article}
% =====================================================================
%  Gemeinsame Vorlage der Arbeitsblätter
%  "Objektorientierte Programmierung mit Java" – Informatik 10 (NTG)
%  Einbinden in jedem Blatt direkt nach \documentclass: % =====================================================================
%  Gemeinsame Vorlage der Arbeitsblätter
%  "Objektorientierte Programmierung mit Java" – Informatik 10 (NTG)
%  Einbinden in jedem Blatt direkt nach \documentclass: % =====================================================================
%  Gemeinsame Vorlage der Arbeitsblätter
%  "Objektorientierte Programmierung mit Java" – Informatik 10 (NTG)
%  Einbinden in jedem Blatt direkt nach \documentclass: \input{ab-vorlage}
%  Übersetzen mit XeLaTeX, LuaLaTeX, pdfLaTeX oder Tectonic.
% =====================================================================
\usepackage[a4paper,left=18mm,right=18mm,top=26mm,bottom=17mm,headheight=15mm,headsep=4mm,footskip=8mm]{geometry}
\usepackage{iftex}
\ifPDFTeX
  \usepackage[T1]{fontenc}
  \usepackage[utf8]{inputenc}
  \usepackage{helvet}
  \usepackage[scaled=0.86]{beramono}
  \renewcommand{\familydefault}{\sfdefault}
\else
  \usepackage{fontspec}
  \setmainfont{texgyreheros}[Extension=.otf,UprightFont=*-regular,BoldFont=*-bold,ItalicFont=*-italic,BoldItalicFont=*-bolditalic]
  \setmonofont{DejaVuSansMono}[Extension=.ttf,UprightFont=*,BoldFont=*-Bold,ItalicFont=*-Oblique,BoldItalicFont=*-BoldOblique,Scale=0.86]
\fi
\usepackage[ngerman,shorthands=off]{babel}
\usepackage{xcolor,graphicx,tabularx,array,enumitem,fancyhdr,tikz,listings,multicol,calc,etoolbox}
\usepackage[most]{tcolorbox}
\usetikzlibrary{arrows.meta,positioning,calc}
\usepackage[hidelinks]{hyperref}
\setlength{\parindent}{0pt}
\setlength{\parskip}{3pt}
\setlist{nosep,leftmargin=*}

% ---------- Kopf- und Fußzeile (einheitliche Form) ----------
\newcommand{\lehrkraft}{Hau}
\newcommand{\abnummer}{}
\pagestyle{fancy}\fancyhf{}
\renewcommand{\headrulewidth}{0.4pt}
\fancyhead[L]{\small Klasse 10\\Datum:}
\fancyhead[C]{\small Themenbereich:\\2. Objektorientierte Programmierung}
\fancyhead[R]{\small \lehrkraft}
\fancyfoot[C]{\scriptsize edu-mrh.de\,·\,Informatik 10\,·\,Blatt \abnummer\,·\,Seite \thepage}
\newcommand{\abtitel}[2]{\renewcommand{\abnummer}{#1}\begin{center}\Large\bfseries #1\quad\underline{#2}\end{center}\vspace{-1mm}}

% ---------- Boxen ----------
\newenvironment{aufgabe}[2][]{\begin{tcolorbox}[enhanced,breakable,colback=white,colframe=black,boxrule=0.7pt,arc=3mm,left=2mm,right=2mm,top=1.5mm,bottom=1.5mm,before skip=2.5mm,after skip=2.5mm]\textbf{Aufgabe #2)}\ifstrempty{#1}{}{ \textit{#1}}\par\vspace{0.6mm}}{\end{tcolorbox}}
\newtcolorbox{merke}[1][Merke]{enhanced,breakable,colback=green!5,colframe=green!40!black,boxrule=0.8pt,arc=2mm,left=2mm,right=2mm,top=1.5mm,bottom=1.5mm,before skip=2.5mm,after skip=2.5mm,title={#1},fonttitle=\bfseries,coltitle=white}
\newtcolorbox{hinweis}{enhanced,breakable,colback=yellow!12,colframe=orange!70!black,boxrule=0.6pt,arc=2mm,left=2mm,right=2mm,top=1mm,bottom=1mm,before skip=2mm,after skip=2mm}
\newcommand{\web}[1]{{\small\color{gray}\textit{Website: #1}}}

% ---------- Lücken, Linien, Karos ----------
\newcommand{\luecke}[1][3.5cm]{\rule[-0.5ex]{#1}{0.4pt}}
\newcommand{\linien}[1]{\par\vspace{1.5mm}\foreach \i in {1,...,#1}{\noindent\rule[0pt]{\linewidth}{0.3pt}\par\vspace{5.5mm}}\vspace{-3mm}}
\newcommand{\kariert}[1]{\par\vspace{1.5mm}\noindent\begin{tikzpicture}\draw[step=5mm,gray!55,thin] (0,0) grid (\linewidth-1pt,#1*5mm);\end{tikzpicture}\par}
\newcommand{\karobox}[2]{\begin{tikzpicture}\draw[step=5mm,gray!55,thin] (0,0) grid (#1,#2);\end{tikzpicture}}
\newcommand{\klassenkarte}[3]{\begin{tabular}{|l|}\hline\textbf{#1}\\\hline #2\\\hline #3\\\hline\end{tabular}}
\newcommand{\code}[1]{\texttt{#1}}

% ---------- Java-Listings ----------
\lstset{language=Java,basicstyle=\ttfamily\small,keywordstyle=\bfseries,commentstyle=\color{gray!80!black},stringstyle=\color{black},showstringspaces=false,columns=flexible,keepspaces=true,tabsize=3,frame=single,framerule=0.4pt,rulecolor=\color{black},xleftmargin=2mm,framexleftmargin=1mm,aboveskip=2mm,belowskip=2mm,breaklines=true,
 literate={ä}{{\"a}}1 {ö}{{\"o}}1 {ü}{{\"u}}1 {Ä}{{\"A}}1 {Ö}{{\"O}}1 {Ü}{{\"U}}1 {ß}{{\ss}}1 {–}{{--}}1 {„}{{\glqq}}1 {“}{{\grqq}}1}
\lstnewenvironment{java}[1][]{\lstset{#1}}{}

%  Übersetzen mit XeLaTeX, LuaLaTeX, pdfLaTeX oder Tectonic.
% =====================================================================
\usepackage[a4paper,left=18mm,right=18mm,top=26mm,bottom=17mm,headheight=15mm,headsep=4mm,footskip=8mm]{geometry}
\usepackage{iftex}
\ifPDFTeX
  \usepackage[T1]{fontenc}
  \usepackage[utf8]{inputenc}
  \usepackage{helvet}
  \usepackage[scaled=0.86]{beramono}
  \renewcommand{\familydefault}{\sfdefault}
\else
  \usepackage{fontspec}
  \setmainfont{texgyreheros}[Extension=.otf,UprightFont=*-regular,BoldFont=*-bold,ItalicFont=*-italic,BoldItalicFont=*-bolditalic]
  \setmonofont{DejaVuSansMono}[Extension=.ttf,UprightFont=*,BoldFont=*-Bold,ItalicFont=*-Oblique,BoldItalicFont=*-BoldOblique,Scale=0.86]
\fi
\usepackage[ngerman,shorthands=off]{babel}
\usepackage{xcolor,graphicx,tabularx,array,enumitem,fancyhdr,tikz,listings,multicol,calc,etoolbox}
\usepackage[most]{tcolorbox}
\usetikzlibrary{arrows.meta,positioning,calc}
\usepackage[hidelinks]{hyperref}
\setlength{\parindent}{0pt}
\setlength{\parskip}{3pt}
\setlist{nosep,leftmargin=*}

% ---------- Kopf- und Fußzeile (einheitliche Form) ----------
\newcommand{\lehrkraft}{Hau}
\newcommand{\abnummer}{}
\pagestyle{fancy}\fancyhf{}
\renewcommand{\headrulewidth}{0.4pt}
\fancyhead[L]{\small Klasse 10\\Datum:}
\fancyhead[C]{\small Themenbereich:\\2. Objektorientierte Programmierung}
\fancyhead[R]{\small \lehrkraft}
\fancyfoot[C]{\scriptsize edu-mrh.de\,·\,Informatik 10\,·\,Blatt \abnummer\,·\,Seite \thepage}
\newcommand{\abtitel}[2]{\renewcommand{\abnummer}{#1}\begin{center}\Large\bfseries #1\quad\underline{#2}\end{center}\vspace{-1mm}}

% ---------- Boxen ----------
\newenvironment{aufgabe}[2][]{\begin{tcolorbox}[enhanced,breakable,colback=white,colframe=black,boxrule=0.7pt,arc=3mm,left=2mm,right=2mm,top=1.5mm,bottom=1.5mm,before skip=2.5mm,after skip=2.5mm]\textbf{Aufgabe #2)}\ifstrempty{#1}{}{ \textit{#1}}\par\vspace{0.6mm}}{\end{tcolorbox}}
\newtcolorbox{merke}[1][Merke]{enhanced,breakable,colback=green!5,colframe=green!40!black,boxrule=0.8pt,arc=2mm,left=2mm,right=2mm,top=1.5mm,bottom=1.5mm,before skip=2.5mm,after skip=2.5mm,title={#1},fonttitle=\bfseries,coltitle=white}
\newtcolorbox{hinweis}{enhanced,breakable,colback=yellow!12,colframe=orange!70!black,boxrule=0.6pt,arc=2mm,left=2mm,right=2mm,top=1mm,bottom=1mm,before skip=2mm,after skip=2mm}
\newcommand{\web}[1]{{\small\color{gray}\textit{Website: #1}}}

% ---------- Lücken, Linien, Karos ----------
\newcommand{\luecke}[1][3.5cm]{\rule[-0.5ex]{#1}{0.4pt}}
\newcommand{\linien}[1]{\par\vspace{1.5mm}\foreach \i in {1,...,#1}{\noindent\rule[0pt]{\linewidth}{0.3pt}\par\vspace{5.5mm}}\vspace{-3mm}}
\newcommand{\kariert}[1]{\par\vspace{1.5mm}\noindent\begin{tikzpicture}\draw[step=5mm,gray!55,thin] (0,0) grid (\linewidth-1pt,#1*5mm);\end{tikzpicture}\par}
\newcommand{\karobox}[2]{\begin{tikzpicture}\draw[step=5mm,gray!55,thin] (0,0) grid (#1,#2);\end{tikzpicture}}
\newcommand{\klassenkarte}[3]{\begin{tabular}{|l|}\hline\textbf{#1}\\\hline #2\\\hline #3\\\hline\end{tabular}}
\newcommand{\code}[1]{\texttt{#1}}

% ---------- Java-Listings ----------
\lstset{language=Java,basicstyle=\ttfamily\small,keywordstyle=\bfseries,commentstyle=\color{gray!80!black},stringstyle=\color{black},showstringspaces=false,columns=flexible,keepspaces=true,tabsize=3,frame=single,framerule=0.4pt,rulecolor=\color{black},xleftmargin=2mm,framexleftmargin=1mm,aboveskip=2mm,belowskip=2mm,breaklines=true,
 literate={ä}{{\"a}}1 {ö}{{\"o}}1 {ü}{{\"u}}1 {Ä}{{\"A}}1 {Ö}{{\"O}}1 {Ü}{{\"U}}1 {ß}{{\ss}}1 {–}{{--}}1 {„}{{\glqq}}1 {“}{{\grqq}}1}
\lstnewenvironment{java}[1][]{\lstset{#1}}{}

%  Übersetzen mit XeLaTeX, LuaLaTeX, pdfLaTeX oder Tectonic.
% =====================================================================
\usepackage[a4paper,left=18mm,right=18mm,top=26mm,bottom=17mm,headheight=15mm,headsep=4mm,footskip=8mm]{geometry}
\usepackage{iftex}
\ifPDFTeX
  \usepackage[T1]{fontenc}
  \usepackage[utf8]{inputenc}
  \usepackage{helvet}
  \usepackage[scaled=0.86]{beramono}
  \renewcommand{\familydefault}{\sfdefault}
\else
  \usepackage{fontspec}
  \setmainfont{texgyreheros}[Extension=.otf,UprightFont=*-regular,BoldFont=*-bold,ItalicFont=*-italic,BoldItalicFont=*-bolditalic]
  \setmonofont{DejaVuSansMono}[Extension=.ttf,UprightFont=*,BoldFont=*-Bold,ItalicFont=*-Oblique,BoldItalicFont=*-BoldOblique,Scale=0.86]
\fi
\usepackage[ngerman,shorthands=off]{babel}
\usepackage{xcolor,graphicx,tabularx,array,enumitem,fancyhdr,tikz,listings,multicol,calc,etoolbox}
\usepackage[most]{tcolorbox}
\usetikzlibrary{arrows.meta,positioning,calc}
\usepackage[hidelinks]{hyperref}
\setlength{\parindent}{0pt}
\setlength{\parskip}{3pt}
\setlist{nosep,leftmargin=*}

% ---------- Kopf- und Fußzeile (einheitliche Form) ----------
\newcommand{\lehrkraft}{Hau}
\newcommand{\abnummer}{}
\pagestyle{fancy}\fancyhf{}
\renewcommand{\headrulewidth}{0.4pt}
\fancyhead[L]{\small Klasse 10\\Datum:}
\fancyhead[C]{\small Themenbereich:\\2. Objektorientierte Programmierung}
\fancyhead[R]{\small \lehrkraft}
\fancyfoot[C]{\scriptsize edu-mrh.de\,·\,Informatik 10\,·\,Blatt \abnummer\,·\,Seite \thepage}
\newcommand{\abtitel}[2]{\renewcommand{\abnummer}{#1}\begin{center}\Large\bfseries #1\quad\underline{#2}\end{center}\vspace{-1mm}}

% ---------- Boxen ----------
\newenvironment{aufgabe}[2][]{\begin{tcolorbox}[enhanced,breakable,colback=white,colframe=black,boxrule=0.7pt,arc=3mm,left=2mm,right=2mm,top=1.5mm,bottom=1.5mm,before skip=2.5mm,after skip=2.5mm]\textbf{Aufgabe #2)}\ifstrempty{#1}{}{ \textit{#1}}\par\vspace{0.6mm}}{\end{tcolorbox}}
\newtcolorbox{merke}[1][Merke]{enhanced,breakable,colback=green!5,colframe=green!40!black,boxrule=0.8pt,arc=2mm,left=2mm,right=2mm,top=1.5mm,bottom=1.5mm,before skip=2.5mm,after skip=2.5mm,title={#1},fonttitle=\bfseries,coltitle=white}
\newtcolorbox{hinweis}{enhanced,breakable,colback=yellow!12,colframe=orange!70!black,boxrule=0.6pt,arc=2mm,left=2mm,right=2mm,top=1mm,bottom=1mm,before skip=2mm,after skip=2mm}
\newcommand{\web}[1]{{\small\color{gray}\textit{Website: #1}}}

% ---------- Lücken, Linien, Karos ----------
\newcommand{\luecke}[1][3.5cm]{\rule[-0.5ex]{#1}{0.4pt}}
\newcommand{\linien}[1]{\par\vspace{1.5mm}\foreach \i in {1,...,#1}{\noindent\rule[0pt]{\linewidth}{0.3pt}\par\vspace{5.5mm}}\vspace{-3mm}}
\newcommand{\kariert}[1]{\par\vspace{1.5mm}\noindent\begin{tikzpicture}\draw[step=5mm,gray!55,thin] (0,0) grid (\linewidth-1pt,#1*5mm);\end{tikzpicture}\par}
\newcommand{\karobox}[2]{\begin{tikzpicture}\draw[step=5mm,gray!55,thin] (0,0) grid (#1,#2);\end{tikzpicture}}
\newcommand{\klassenkarte}[3]{\begin{tabular}{|l|}\hline\textbf{#1}\\\hline #2\\\hline #3\\\hline\end{tabular}}
\newcommand{\code}[1]{\texttt{#1}}

% ---------- Java-Listings ----------
\lstset{language=Java,basicstyle=\ttfamily\small,keywordstyle=\bfseries,commentstyle=\color{gray!80!black},stringstyle=\color{black},showstringspaces=false,columns=flexible,keepspaces=true,tabsize=3,frame=single,framerule=0.4pt,rulecolor=\color{black},xleftmargin=2mm,framexleftmargin=1mm,aboveskip=2mm,belowskip=2mm,breaklines=true,
 literate={ä}{{\"a}}1 {ö}{{\"o}}1 {ü}{{\"u}}1 {Ä}{{\"A}}1 {Ö}{{\"O}}1 {Ü}{{\"U}}1 {ß}{{\ss}}1 {–}{{--}}1 {„}{{\glqq}}1 {“}{{\grqq}}1}
\lstnewenvironment{java}[1][]{\lstset{#1}}{}

\renewcommand{\themenbereich}{Grundwissen 6 bis 9}
\begin{document}
\abtitel{0}{Grundwissen}

\textbf{0.1 Objektorientierte Modellierung -- Klassen und Objekte}

\begin{center}
\begin{tikzpicture}[x=1cm,y=1cm,font=\small]
\node[anchor=north west] (kk) at (0,0) {\klassenkarte{Dreieck}{seitenlaenge\\farbe}{setzeFarbe(neu)}};
\node[anchor=south west] at (0,0.15) {\luecke[5cm]};
\node[anchor=north west,draw,rounded corners=3mm,inner sep=0pt] (ok) at (8.5,0) {\begin{tabular}{l}\rule{0pt}{4mm}\textbf{form1: Dreieck}\\[1mm]\hline\rule{0pt}{4mm}seitenlaenge = 20\\farbe = "blau"\\[1mm]\end{tabular}};
\node[anchor=south west] at (8.5,0.15) {\luecke[5cm]};
% Klammern Klassenkarte
\draw[decorate,decoration={brace,amplitude=4pt}] ([xshift=2pt]kk.north east) -- ([xshift=2pt]kk.north east|-{(0,-0.55)}) node[midway,right=6pt] {\luecke[3cm]};
\draw[decorate,decoration={brace,amplitude=4pt}] ([xshift=2pt]kk.north east|-{(0,-0.65)}) -- ([xshift=2pt]kk.north east|-{(0,-1.55)}) node[midway,right=6pt] {\luecke[3cm]};
\draw[decorate,decoration={brace,amplitude=4pt}] ([xshift=2pt]kk.north east|-{(0,-1.65)}) -- ([xshift=2pt]kk.south east) node[midway,right=6pt] {\luecke[3cm]};
% Klammern Objektkarte
\draw[decorate,decoration={brace,amplitude=4pt}] ([xshift=2pt]ok.north east) -- ([xshift=2pt]ok.north east|-{(0,-0.6)}) node[midway,right=6pt] {\luecke[3cm]};
\draw[decorate,decoration={brace,amplitude=4pt}] ([xshift=2pt]ok.north east|-{(0,-0.7)}) -- ([xshift=2pt]ok.south east) node[midway,right=6pt] {\luecke[3cm]};
\end{tikzpicture}
\end{center}

\textbf{0.2 Algorithmik}\quad \textbf{0.2.1 Ablaufbeschreibung durch Algorithmen}
\begin{merke}[Definition]
Ein Algorithmus ist eine \luecke[3.2cm] Abfolge von \luecke[3.2cm] und \luecke[3.5cm] Anweisungen zur Lösung eines Problems. Die Anweisungen werden in einer festgelegten Reihenfolge durchlaufen und führen von einem Anfangszustand (Aufgabenstellung) zu einem konkreten Endzustand (Lösung).

Kann ein Algorithmus auch maschinell ausgeführt werden, so nennt man ihn ein \luecke[3.2cm].
\end{merke}

\textbf{0.2.2 Bausteine von Algorithmen -- Struktogramme}

\begin{minipage}[t]{0.47\linewidth}
\textbf{Anweisung}\\[1mm]
\begin{tikzpicture}[x=1cm,y=1cm]\draw (0,0) rectangle (4,0.7); \node[anchor=west] at (0.15,0.35) {schritt()};\end{tikzpicture}
\end{minipage}\hfill
\begin{minipage}[t]{0.47\linewidth}
\textbf{Sequenz}\\[1mm]
\begin{tikzpicture}[x=1cm,y=1cm,font=\ttfamily\small]
\draw (0,0) rectangle (4.5,2.1); \draw (0,0.7) -- (4.5,0.7); \draw (0,1.4) -- (4.5,1.4);
\node[anchor=west] at (0.15,1.75) {linksDrehen()}; \node[anchor=west] at (0.15,1.05) {schritt()}; \node[anchor=west] at (0.15,0.35) {aufheben()};
\end{tikzpicture}
\end{minipage}

\vspace{2mm}
\textbf{Die bedingte Anweisung}

\begin{minipage}[t]{0.47\linewidth}
\textit{einseitig}\\[1mm]
\begin{tikzpicture}[x=1cm,y=1cm]
\draw (0,0) rectangle (5.5,2.2); \draw (0,2.2) -- (2.75,1.2) -- (5.5,2.2); \draw (0,1.2) -- (5.5,1.2); \draw (2.75,0) -- (2.75,1.2); \draw (2.75,1.2) -- (5.5,0);
\end{tikzpicture}\\[2mm]
\textbf{Java:}
\begin{java}
if (istZiegel()) {
   aufheben();
   schritt();
}
\end{java}
\end{minipage}\hfill
\begin{minipage}[t]{0.47\linewidth}
\textit{zweiseitig}\\[1mm]
\begin{tikzpicture}[x=1cm,y=1cm,font=\small]
\draw (0,0) rectangle (5.5,2.2); \draw (0,2.2) -- (2.75,1.2) -- (5.5,2.2); \draw (0,1.2) -- (5.5,1.2); \draw (2.75,0) -- (2.75,1.2);
\node at (2.75,1.85) {\ttfamily\itshape istWand()}; \node at (0.6,1.4) {wahr}; \node at (4.9,1.4) {falsch};
\node at (1.35,0.6) {\ttfamily linksDrehen()}; \node at (4.1,0.6) {\ttfamily schritt()};
\end{tikzpicture}\\[2mm]
\textbf{Java:}\\[1mm]
\framebox[\linewidth]{\begin{minipage}{0.93\linewidth}\vspace{1mm}\linien{4}\vspace{-2mm}\end{minipage}}
\end{minipage}

\newpage
\textbf{Die Wiederholung}

\begin{minipage}[t]{0.47\linewidth}
\textit{bedingte Wiederholung}\\[1mm]
\begin{tikzpicture}[x=1cm,y=1cm,font=\small]
\draw (0,0) rectangle (5.5,1.6); \draw (0.7,0) -- (0.7,1.0) -- (5.5,1.0);
\node[anchor=west] at (0.15,1.3) {Wiederhole solange {\ttfamily\itshape nichtIstWand()}}; \node[anchor=west] at (0.9,0.5) {\ttfamily schritt()};
\end{tikzpicture}\\[2mm]
\textbf{Java:}\\[1mm]
\framebox[\linewidth]{\begin{minipage}{0.93\linewidth}\vspace{1mm}\linien{2}\vspace{-2mm}\end{minipage}}
\end{minipage}\hfill
\begin{minipage}[t]{0.47\linewidth}
\textit{Wiederholung mit fester Anzahl}\\[1mm]
\begin{tikzpicture}[x=1cm,y=1cm,font=\small]
\draw (0,0) rectangle (5.5,1.6); \draw (0.7,0) -- (0.7,1.0) -- (5.5,1.0);
\node[anchor=west] at (0.15,1.3) {Wiederhole \luecke[3cm]};
\end{tikzpicture}\\[2mm]
\textbf{Java:}
\begin{java}
for (int i = 0; i < 5; i++) {
   hinlegen();
   schritt();
}
\end{java}
\end{minipage}

\textbf{0.3 Vererbung}

Klassen können voneinander \textbf{erben}. Attribute und Methoden der Oberklasse sind dann auch in Unterklassen verfügbar. \underline{Wichtig}: Eine Klasse kann in Java immer nur von \underline{einer} Klasse erben (keine Mehrfachvererbung). Eine hierarchische Vererbung ist jedoch möglich.

\begin{aufgabe}{1}
\textbf{a)} Ordne die Klassenkarten zu einer sinnvollen Vererbungshierarchie (Pfeile einzeichnen, oder rechts neu zeichnen).\\
\textbf{b)} Ein Hai hat damit die Attribute und Methoden:
\end{aufgabe}
\begin{minipage}[t]{0.55\linewidth}
\vspace{0pt}
\begin{tikzpicture}[x=1cm,y=1cm,font=\small]
\node[anchor=north west] at (5.2,0) {\klassenkarte{Tier}{alter\\gewicht}{fressen()}};
\node[anchor=north west] at (2.2,-1.15) {\klassenkarte{Fisch}{\rule{0pt}{4mm}}{schwimmen()}};
\node[anchor=north west] at (0,-2.55) {\klassenkarte{Clownfisch}{\rule{0pt}{4mm}}{nemoSuchen()}};
\node[anchor=north west] at (4.4,-2.95) {\klassenkarte{Hai}{anzahlZaehne}{\rule{0pt}{4mm}}};
\end{tikzpicture}
\end{minipage}\hfill
\begin{minipage}[t]{0.42\linewidth}
\vspace{0pt}\linien{4}
\end{minipage}

\textbf{0.4 Datenbanken}\quad \textbf{0.4.1 Von der Klasse zur Datenbank}

\begin{minipage}[t]{0.28\linewidth}
\vspace{0pt}
\klassenkarte{Mitglied}{name\\vorname\\\ldots\\beitrag\\gezahlt}{\ldots}
\end{minipage}\hfill
\begin{minipage}[t]{0.68\linewidth}
\vspace{0pt}\scriptsize
\begin{tabular}{|l|l|l|l|l|r|l|}
\multicolumn{7}{c}{\textbf{Mitglied}}\\\hline
\underline{Name} & Vorname & Adresse & Eintritt & Sportart & Beitrag & gezahlt\\\hline
Hopf & Hans & Bahnstr.\,\ldots & 2020-02-01 & Volleyball & 90 & WAHR\\
Paulus & Petra & Opernpl.\,\ldots & 2018-10-01 & Tennis & 90 & FALSCH\\
Leroy & Jana & An der\,\ldots & 2016-08-01 & Handball & 45 & WAHR\\\hline
\end{tabular}\\[1mm]
\normalsize Klassenname $\rightarrow$ \luecke[2.4cm]\quad Attributname $\rightarrow$ \luecke[2.4cm]\\[1mm] Objekt $\rightarrow$ \luecke[2.4cm]\quad Primärschlüssel $\rightarrow$ \luecke[2.4cm]
\end{minipage}

\begin{merke}\small
Relationale Datenbanken speichern Daten in \textbf{Tabellen}. Das \textbf{Tabellenschema} beschreibt den Aufbau der Tabelle durch Angabe der \textbf{Attributnamen} und deren \textbf{Datentypen}. Jeder \textbf{Datensatz} entspricht einer Zeile in der Tabelle und wird durch seinen Primärschlüsselwert eindeutig identifiziert. Der \textbf{Primärschlüssel} wird durch Unterstreichung in der Tabelle und im Tabellenschema hervorgehoben.
\end{merke}

\textbf{0.4.2 SQL -- } \luecke[9cm]

\small\begin{tabular}{@{}l l l@{}}
\code{SELECT [DISTINCT]} & \code{<Attribut1>, <Attribut2>, ...} & $\leftarrow$ \luecke[5cm]\\
\code{FROM} & \code{<Tabellenname>} & \\
\code{WHERE} & \code{<Bedingung>} & $\leftarrow$ \luecke[5cm]\\
\code{ORDER BY} & \code{<Attribut> [DESC | ASC]} & $\leftarrow$ \luecke[5cm]\\
\end{tabular}
\end{document}
